\section[Banach--Colmez 空间与向量丛的分类]{Banach--Colmez 空间与向量丛的分类（11 月 23 日）}
% 三级小节 8.1.1--8.1.3、8.3.1--8.3.4 需要编号
\setcounter{secnumdepth}{3}

\srcpage{37}

固定 $E$ 与 $\pi$。对 $S\in\Perf_{\Fq}$，我们简记 $Y_S:=Y_{S,E}$，$X_S:=X_{S,E}$。

\subsection{回顾上次内容}

\subsubsection{$Y$ 上的除子}

上次我们谈的是「$Y_{S,E}$ 到 $S$ 的截面」。我们曾说明，下列集合之间典范地存在双射。
\begin{itemize}
  \item $Y_{S,E}^\diamond\to S$ 的截面。
  \item 映射 $S\to(\Spa E)^\diamond$。
  \item $S$ 的 $E$-untilt $S^\#$。
  \item $Y_{S,E}$ 中的 1 次闭 Cartier 除子 $D$。（用上一条的说法，$D\cong S^\#$。）
\end{itemize}
模问题 $\Div^1_Y$ 分类这些集合。

\subsubsection{$X$ 上的除子}

对 $X_S$ 也有一个变体。下列集合之间典范地存在双射。
\begin{itemize}
  \item 映射 $S\to(\Spa E)^\diamond/\varphi_S^{\mathbb{Z}}$。
  \item $X_S$ 中的 1 次闭 Cartier 除子 $D$。
\end{itemize}
模问题 $\Div^1_X$ 分类这些集合。\footnote{也可以用「模 Frobenius 的 untilt」来描述，但那需要层化。}

我们常在 $\Perf_{\overline{\Fq}}$ 而不是 $\Perf_{\Fq}$ 上工作。此时得到
\[
\Div^1_Y\cong(\Spa\breve E)^\diamond,\qquad
\Div^1_X\cong(\Spa\breve E)^\diamond/\varphi_S^{\mathbb{Z}}.
\]
另一个表示会很有用。取 $C=\widehat{\breve E}$。则
\[
(\Spa\breve E)^\diamond/\varphi_S^{\mathbb{Z}}
=(\Spa C)^\diamond/(\underbrace{I_E\rtimes\varphi_S^{\mathbb{Z}}}_{W_E}),
\]
其中 $W_E$ 是 $E$ 的 Weil 群。因此 $\pi_1(\Div^1_X)\cong W_E$。

\subsubsection{基本群}

将有一个映射 $\Div^1_X\to\Pic^1_X$，其中 $\Pic^1_X$ 是 1 次线丛的模空间，它把 $D$ 送到 $\cO(D)$。线丛 $\cO(1)$ 给出一个特定的映射 $\mathrm{pt}\to\Pic^1_X$。上次的讨论用 Lubin--Tate 群算出了这个纤维积：
\[
\begin{tikzcd}
(\Spa\breve E_\infty)^\diamond \arrow[r] \arrow[d] & \mathrm{pt} \arrow[d]\\
\Div^1_X \arrow[r,"AJ^1"] & \Pic^1_X
\end{tikzcd}
\]
设 $G/\cO_{\breve E}$ 是满足 $\log_G(X)=X+\frac{1}{\pi}X^q+\frac{1}{\pi^2}X^{q^2}+\cdots$ 的 Lubin--Tate 群。于是我们定义了 $\breve E_\infty=\breve E(G[\pi^\infty])^{\wedge}$。

我们曾说明，下列集合之间典范地存在双射。\footnote{不过这些双射依赖于 $\pi$ 的选取。}
\begin{itemize}
  \item 映射 $S\to(\Spa\breve E_\infty)^\diamond=\Spa\breve E_\infty^\flat$。（具体地，$\breve E_\infty^\flat\cong\Fq((X^{1/p^\infty}))$。）
  \item $X_S$ 中的 1 次闭 Cartier 除子 $D$，连同同构 $\cO(D)\cong\cO(1)$。
\end{itemize}
上图中的竖直映射都是 $E^\times$-主丛。

Abel--Jacobi 映射 $AJ^1:\Div^1_X\to\Pic^1_X=[\mathrm{pt}/E^\times]$ 给出映射 $W_E=\pi_1(\Div^1_X)\to\pi_1(\Pic^1_X)=E^\times$，它就是局部类域论的 Artin 互反律映射。

Fargues 在 [F17] 中以此为基础，模仿 Deligne 对几何类域论的论证，给出了局部类域论的一个新证明。

\srcpage{38}

（$W_E$ 的任意一维特征标都通过下降 $d$ 次外幂，在 $\Div^d_X=(\Div^1_X)^d/\Sigma_d$ 上诱导出秩 1 局部系统。关键在于证明 $AJ^d$ 的纤维对 $d\gg0$ 单连通，从而可以下降到 $\Pic^d_X$。）

\subsection{Banach--Colmez 空间}

参考文献：[LB18]。

设 $S\in\Perf_{\Fq}$，$\mathcal{E}$ 是 $X_S$ 上的向量丛（即有限秩局部自由 $\cO_{X_S}$-模）。

Kedlaya--Liu 的结果表明，adic 空间上向量丛的概念性质良好。

\begin{theorem}[Kedlaya--Liu]
设 $X=\Spa(A,A^+)$ 是 affinoid 解析 adic 空间（故 $\cO_X$ 是层）。则我们有范畴等价
\[
\mathrm{VB}(X)\xleftarrow{\ \sim\ }\{\text{有限投射 }A\text{-模}\}/\sim,
\]
它把 $M\otimes_A\cO_X$ 送到 $M$。
此外，对 $\mathcal{E}\in\mathrm{VB}(X)$，$i>0$ 时 $H^i(X,\mathcal{E})=0$；若 $\cO_X$ 是 étale 层，对 $H^i_{\text{ét}}$ 也有同样的结论。\footnote{对我们来说这总是成立的。通常那些表明 $\cO_X$ 在解析拓扑下是层的判别法，也会表明它在 étale 拓扑下是层。}
\end{theorem}

\begin{proposition}
若 $S$ 是 affinoid，则 $i\ge2$ 时 $H^i(X_S,\mathcal{E})=0$，$i\ge1$ 时 $H^i(Y_S,\mathcal{E})=0$。
\end{proposition}

\begin{proof}[证明梗概]
取 $\varpi$ 为伪一致化子。我们有半径函数 $\mathrm{rad}\colon Y_S\to(0,\infty)$，用它来比较 $|[\varpi]|$ 与 $|\pi|$。在我们的正规化下，$\varphi_S$ 把半径乘以 $q$。

对区间 $J=[a,b]$（$a,b\in\mathbb{Q}$），有有理子集 $Y_{S,J}\subset Y_S\subset\Spa W_{\cO_E}(R^+)$。它可以描述为
\[
\{|[\varpi]|^b\le|\pi|\le|[\varpi]|^a\neq0\}=Y_{S,J}\subset\mathrm{rad}^{-1}(J).
\]
我们指出，这个包含在秩 1 点上取等号；但 $Y_{S,J}$ 是开的，而 $\mathrm{rad}^{-1}(J)$ 是闭的（因为 $J$ 是闭区间）。于是 $Y_{S,J}$ 是 affinoid 且解析的。此时 $X_S$ 也可以表成
\[
Y_{S,[1,q]}\big/\bigl(Y_{S,[1,1]}\xrightarrow{\sim_\varphi}Y_{S,[q,q]}\bigr).
\]
这个覆盖允许构造计算 $X_S$ 上同调的 \v{C}ech 复形：
\[
R\Gamma(X_S,\mathcal{E})\cong
\Bigl[\mathcal{E}(Y_{S,[1,q]})\xrightarrow{\varphi^{-1}}\mathcal{E}(Y_{S,[q,q]})\Bigr].
\]
这给出次数 $\ge2$ 的消没。

对 $Y_S$，我们写 $Y_S=\bigcup_J Y_{S,J}$，且过渡映射 $\cO(Y_{S,J'})\to\cO(Y_{S,J})$ 的像稠密。于是
\[
R\Gamma(Y_S,\mathcal{E})=R\varprojlim_{I}\mathcal{E}(Y_{S,I}),
\]
而 $R^1\varprojlim$ 由 Mittag--Leffler 性质消没。口号是：「$Y_S$ 是 Stein 的」。
\end{proof}

\begin{remark}
我们本也可以由 $Y$ 的消没，用群上同调推出 $X$ 的消没。然而，上述证明还给出了用 affinoid 计算 $X$ 上同调的配方，这在实践中很有用。
\end{remark}

\srcpage{39}

\begin{proposition}
函子 $T\in\Perf_{/S}\mapsto H^0(X_T,\mathcal{E}|_{X_T})$ 是 v-层。事实上，$T\mapsto R\Gamma(X_T,\mathcal{E}|_{X_T})$ 是 $D(\mathbb{Z})$ 中复形的 v-层。

特别地，若对一切 $T\in\Perf_{/S}$ 都有 $H^0(X_T,\mathcal{E}|_{X_T})=0$，则 $T\mapsto H^1(X_T,\mathcal{E}|_{X_T})$ 是 v-层。
\end{proposition}

\begin{proof}
只需在 $\widehat{\otimes}_E E_\infty$ 之后验证，因为 $\mathcal{E}\hookrightarrow E_\infty$ 作为 Banach $E$-向量空间的包含分裂（例如由 Hahn--Banach）。要点是 $E_\infty$ 是完满的，且 $X_S\times_{\Spa E}\Spa E_\infty$ 是完满空间。因此 $S$ 上的 v-覆盖诱导出这个对象上的 v-覆盖。现在我们对一般的完满空间使用 v-层与无环性性质。
\end{proof}

\begin{definition}
(1) 函子 $BC(\mathcal{E})\colon\Perf_{/S}\to\mathrm{Sets}$，$T\mapsto H^0(X_T,\mathcal{E}|_{X_T})$，称为 $\mathcal{E}$ 对应的 \emph{Banach--Colmez 空间}。（对通常的曲线，这会是一个维数为 $h^0(X,\mathcal{E})$ 的仿射空间。）

(2) 若 $BC(\mathcal{E})=0$，则定义 $BC(\mathcal{E}[1])\colon T\to H^1(X_T,\mathcal{E}|_{X_T})$ 为 $\mathcal{E}$ 对应的\emph{负 Banach--Colmez 空间}。
\end{definition}

\begin{proposition}
(1) $BC(\mathcal{E})$、$BC(\mathcal{E}[1])$ 都是 $S$ 上的局部空间菱形。

(2) 若 $E=\Fq((t))$，则 $BC(\mathcal{E})$ 甚至可由完满空间表示。

(3) 若 $\mathcal{E}=\cO(\lambda)$，其中 $0<\lambda\le[E:\mathbb{Q}_p]$，写 $\lambda=\frac{r}{s}$ 且 $(s,r)=1$、$r,s>0$，则 $BC(\mathcal{E})\cong\widetilde{\mathbb{D}}^r_S$ 是 $S$ 上的 $r$ 维完满圆盘；且当 $S$ 为 affinoid 时 $H^1(X_S,\mathcal{E})=0$。

(4) $R\Gamma(X_S,\cO_{X_S})\cong R\Gamma_{\text{proét}}(S,\underline{E})$。特别地，$R\Gamma(X_C,\cO_{X_C})\cong E[0]$。
\end{proposition}

\begin{proof}[证明梗概]
从 (3) 开始。这与上次的计算类似，当时我们算出了 $BC(\cO(1))\cong\widetilde{\mathbb{D}}_S$。

为简单起见设 $E=\mathbb{Q}_p$。则 $BC(\cO(\lambda))\cong\widetilde{G}_S$，其中 $\widetilde{G}$ 是 $p$ 可除群 $G/\Fq$ 的万有覆盖，后者的 Dieudonné 模为 $D_{-\lambda}$。

$H^1$ 的消没可以用上述 \v{C}ech 复形直接计算证明。细节留作练习。

(4) 在 $S$ 上 pro-étale 局部，我们使用正合列
\[
0\to\cO_{X_S}\to\cO_{X_S}(1)\to\cO_{S^\#}\to0.
\]
取整体截面，得到
\[
\begin{tikzcd}[column sep=small]
0 \arrow[r] & H^0(\cO) \arrow[r] & H^0(\cO(1)) \arrow[r,"\log_{\widetilde G}"] \arrow[d,equal] & H^0(\cO_{S^\#}) \arrow[r] & H^1(\cO) \arrow[r] & 0\\
& & \widetilde{G}(S^\#) & & &
\end{tikzcd}
\]
现在利用 $\log_{\widetilde G}$ 是 pro-étale 局部满射，其核等同于 $E$。\yizhu{原文此处把第三项印作 $R^{\#}$，疑为 $H^0(\cO_{S^\#})$ 之误；译文按正合列的含义写出。}{A（原文 §8.2 p.39）}

(1) 与 (2)：用各种正合列从 (3) 自举。例如 $BC(\cO(-1)[1])\cong(\mathbb{A}^1_E)^\diamond/\underline{E}$：对 $S/(\Spa E^{\mathrm{LT}}_\infty)^\diamond$ 使用
\[
0\to\cO_{X_S}(-1)\to\cO_{X_S}\to\cO_{S^\#}\to0,
\]
便得到
\[
0\to H^0(\cO)\to H^0(\cO_{S^\#})\to H^1(\cO(-1))\to H^1(\cO),
\]
并利用 $H^0(\cO_{S^\#})=(\mathbb{A}^1_E)^\diamond(S)$、$H^0(\cO)=\underline{E}(S)$、$H^1(\cO)$ 局部为零。这就处理了平移 1 的情形。$\cO(1/n)$ 的情形可以直接处理。
\end{proof}

\srcpage{40}

\subsection{向量丛的分类}

回到 $S=\Spa C$ 是几何点的情形。

\begin{theorem}
我们有双射
\[
\Isoc_E/\sim\ \longleftrightarrow\ \mathrm{VB}(X_C)/\sim,
\]
具体地说，任意 $\mathcal{E}\in\mathrm{VB}(X_C)$ 都同构于
\[
\bigoplus_{\lambda\in\mathbb{Q}}\cO_{X_C}(\lambda)^{n_\lambda},
\]
其中 $n_\lambda\in\mathbb{Z}$ 唯一。
\end{theorem}

\subsubsection{第一步}

先确立 $\cO_{X_C}(1)$ 是丰沛的。这意味着对任意 $\mathcal{E}$ 及一切 $n\gg0$，$\mathcal{E}(n)$ 由整体截面生成，且 $H^1(X_C,\mathcal{E}(n))=0$。（这对任意 affinoid $S$ 都成立：对任意 affinoid $S$，Fargues--Fontaine 曲线都有代数版本。）

这是 Kedlaya--Liu 的定理，用上述 \v{C}ech 复形与显式估计硬算得到。任务是找到足够多的 Frobenius 不变量（模 Frobenius 扭），要点在于扭得足够多能使某些级数收敛。

\subsubsection{第二步}

证明 $\Pic(X_C)\cong\mathbb{Z}$，同构由 $n\mapsto\cO_{X_C}(n)$ 给出。第一步表明，任意 $\mathcal{L}\in\Pic(X_C)$ 在概形曲线上都是生成平凡的。于是 $\mathcal{L}\cong\cO(D)$，其中 $D$ 是概形曲线上的某个除子。但概形曲线上的所有闭点都对应 untilt，而它们的除子都同构于 $\cO(1)$。因此 $\cO(D)\cong\cO(\deg D)$，其中 $\deg$ 按系数之和朴素地定义。这说明 $\mathbb{Z}\twoheadrightarrow\Pic(X_C)$。由于 $n>0$ 时 $H^0(\cO(-n))=0$，没有非平凡关系，所以这个满射是同构。

\subsubsection{第三步}

建立 Harder--Narasimhan 形式化。函数 $\mathrm{rank},\deg\colon\mathrm{VB}(X_C)\to\mathbb{Z}$ 给出斜率函数 $\mu:=\frac{\deg}{\mathrm{rank}}$。于是由形式理由得到 Harder--Narasimhan 滤过。

利用命题 8.6 的 $\lambda\ge0$ 时 $H^1(X_C,\cO_{X_C}(\lambda))=0$，可知 Harder--Narasimhan 滤过分裂，从而化归到 $\mathcal{E}$ 半稳定的情形。再稍加论证可化归到 $\mathcal{E}$ 是斜率 0 的半稳定丛的情形。

\subsubsection{第四步}

目标是证明：任意斜率 0 的半稳定 $\mathcal{E}$ 实际上平凡。

我们断言，只需在可能扩大 $C$ 之后证明这一点。这由 v-下降得到：若它在 $C'/C$ 上成立（利用 $\Spa C'\to\Spa C$ 是 v-覆盖），则同构 $\mathcal{E}\cong\cO^n$ 的 torsor 是 $\Spa C$ 上的 $\mathrm{GL}_n(E)$-v-主丛。由 [S17] 中的一个定理\footnote{看起来是 [S17, Lemma 10.13]。}，任何这样的 v-主丛都来自 pro-étale 主丛，而 $\Spa C$ 上的任意 pro-étale 主丛都分裂。于是可以用 v-下降在底上找到分裂。

此外，我们可以归纳地假设定理在更小的秩上成立。考虑最小的 $d\ge0$（在扩大 $C$ 的意义下），使得存在单射
\[
0\to\cO_{X_C}(-d)\to\mathcal{E}\to\overline{\mathcal{E}}\to0.
\]
若 $d=0$，则由归纳并利用 $H^1(X_C,\cO_{X_C})=0$ 即可。若 $d\ge2$，则是相当简单的矛盾。关键情形是 $d=1$。此时得到
\[
0\to\cO_{X_C}(-1)\to\mathcal{E}\to\overline{\mathcal{E}}\to0.
\]
于是 $\overline{\mathcal{E}}$ 的秩为 $n-1$、次数为 1、斜率 $\ge0$。由归纳必有 $\overline{\mathcal{E}}\cong\cO^i_{X_C}\oplus\cO_{X_C}\bigl(\frac{1}{n-1-i}\bigr)$。关键情形是 $\overline{\mathcal{E}}=\cO\bigl(\frac{1}{n-1}\bigr)$。这化归为下面的引理。

\srcpage{41}

\begin{lemma}
设 $\mathcal{E}$ 是一个扩张
\[
0\to\cO_{X_C}(-1)\to\mathcal{E}\to\cO_{X_C}(1/n)\to0.
\]
则在可能扩大 $C$ 之后，$H^0(X_C,\mathcal{E})\neq0$。\yizhu{原文引理中写作 $1/n$，而上一页末把关键情形写作 $\cO(\frac{1}{n-1})$；两处指标不一致，疑为原文笔误，译文照印。}{A（原文 §8.3 p.40--41）}
\end{lemma}

\begin{remark}
化归到这个引理可以追溯到 [HP04]，而且在已知的所有分类证明中都是一样的。

Hartl--Pink 对该引理的证明（那是等特征情形）是一个困难的计算。Kedlaya--Liu 用更困难的计算把它推广到混合特征。Fargues--Fontaine 用 $p$ 可除群的 Lubin--Tate 与 Drinfeld 模空间上的周期映射，给出了一个漂亮的几何论证。下面是新的证明。
\end{remark}

\begin{proof}
反设不然。则对一切 $S\in\Perf_{/C}$ 都有单射
\[
H^0(X_S,\cO_{X_S}(1/n))\hookrightarrow H^1(X_S,\cO_{X_S}(-1)),
\]
即 $BC(\cO(1/n))\hookrightarrow BC(\cO(-1)[1])$。几何上，这就是
\[
\widetilde{\mathbb{D}}_C\hookrightarrow(\mathbb{A}^1_{C^\#})^\diamond/\underline{E}.
\]
我们断言这个映射必然也是满射。（这是一个例子，说明把菱形看成真正的几何对象、而不是「仅仅」看成函子是有利的。）像不可能包含在经典点中，因为目标中的经典点完全间断，而 $\widetilde{\mathbb{D}}_C$ 连通。所以像包含某个非经典点。由对仿射线的理解，在可能扩大 $C$ 之后，像包含一个非空开子集（非经典点的一般性质：延拓标量后，其纤维包含一个开子集\footnote{例子：在 $\mathbb{A}^1$ 上，非经典点 $x$ 就是「一般点 $B(x,r)$」。基变到剩余域 $C(x)$ 后，有一个重言点 $\tilde x$，我们断言 $B(\tilde x,<r)$ 落在 $\{x\}\subset|(\mathbb{A}^1_C)^{\mathrm{ad}}|$ 的原像中。你可以把 $B(x,r)$ 的 Gauß 点想成经典点周围所有更小球的外部。基变之后，出现了新的经典点，它们的球还没有被删掉。}）。由平移，像包含 0 的一个非空开邻域。于是映射是满射，因为两边都是 $E$-向量空间（所以任意拟紧开子集都可以用 $\pi$ 的幂缩放到 0 的小邻域中），从而它是同构。

但这个映射不可能是同构，因为容易看出 $(\mathbb{A}^1_{C^\#})^\diamond/\underline{E}$ 不可表（它没有整体函数）。
\end{proof}
